\section{Marco teórico}

\subsection{Lógica combinacional y descripción en VHDL}
En un circuito combinacional, las salidas dependen de las entradas
actuales y no de un estado almacenado. En VHDL, la entidad define
los puertos y la arquitectura describe el comportamiento o las
conexiones entre los bloques. Aunque las instrucciones dentro de un
proceso se evalúan en orden, los procesos y las asignaciones
concurrentes representan circuitos que operan en paralelo~\cite{fundamentos}.

Un proceso combinacional debe ser sensible a todas las señales que
lee y asignar sus salidas para cada alternativa de ejecución. Las
asignaciones incompletas pueden inferir latches, mientras que una
lista de sensibilidad incompleta puede producir diferencias entre
la simulación y el hardware sintetizado. Una sentencia \texttt{case}
con valores por defecto y una rama \texttt{when others} permite
describir de forma completa la selección por opcode~\cite{combinacional}.

Las variables, asignadas con \texttt{:=}, actualizan su valor
inmediatamente dentro del proceso. Las señales, asignadas con
\texttt{<=}, programan una actualización posterior. Esta diferencia
resulta relevante al calcular un resultado temporal y derivar de
él una bandera dentro del mismo proceso~\cite{fundamentos}.

\subsection{Tipos numéricos y aritmética de cuatro bits}
El paquete \texttt{ieee.std\_logic\_1164} define los tipos
\texttt{std\_logic} y \texttt{std\_logic\_vector}. El paquete
\texttt{ieee.numeric\_std} incorpora los tipos \texttt{unsigned} y
\texttt{signed}, que hacen explícita la interpretación numérica de
los vectores y permiten aplicar operaciones aritméticas y
comparaciones~\cite{combinacional}. En este laboratorio se utiliza
la interpretación sin signo indicada en la consigna.

Un operando sin signo de cuatro bits representa valores entre
$0$ y $2^4-1=15$. Si se conservan únicamente los cuatro bits menos
significativos de un resultado aritmético, el valor se reduce
módulo $16$. Por ejemplo, $15+1$ produce \texttt{0000} con acarreo,
y $0-1$ produce \texttt{1111} con préstamo. El resultado binario
de la resta sigue siendo un vector sin signo; no representa un
número negativo bajo esta interpretación~\cite{casoalu}.

Para obtener el acarreo de una suma es necesario extender los
operandos antes de operar. Si $T$ es el resultado extendido:
\begin{equation}
T=(0\mathbin{\&}A)+(0\mathbin{\&}B),\qquad
R=T(3\!:\!0),\quad C=T(4),
\end{equation}
donde $\&$ indica concatenación. La extensión conserva el quinto
bit que se perdería al calcular directamente en cuatro bits.

\subsection{Banderas de estado}
Según la especificación, Carry expresa el acarreo de salida en
$A+B$ y el préstamo (\textit{borrow}) en $A-B$. En la resta sin
signo existe préstamo cuando $A<B$. Para las operaciones lógicas
y las comparaciones, la consigna fija Carry en cero~\cite{guialab2}.
El acarreo sin signo y el desbordamiento de una operación con signo
son condiciones diferentes y dependen de la interpretación de los bits.
Para las operaciones restantes se adopta la convención de la
Tabla~\ref{tab:carry}: acarreo al incrementar, préstamo al decrementar
o calcular $0-X$, y bit descartado en los desplazamientos lógicos.
Esta extensión es una decisión de diseño del equipo.

Zero se obtiene del resultado final:
\begin{equation}
Z=\begin{cases}
1,& R=\texttt{0000},\\
0,& R\ne\texttt{0000}.
\end{cases}
\end{equation}
Por tanto, Zero no indica por sí sola si los operandos son iguales.
En la operación de comparación $A=B$, un resultado verdadero se
codifica como \texttt{0001} y deja Zero en cero.

\subsection{Complemento a dos y reutilización del sumador}
Para un vector $X$ de cuatro bits, su complemento a dos corresponde
a $(-X)\bmod16$ y se obtiene invirtiendo sus bits y sumando uno.
La misma identidad permite expresar una resta como una suma:
\begin{equation}
A-B\equiv A+\overline{B}+1\pmod{16}.
\end{equation}
Un sumador con entrada de acarreo permite implementar suma, resta,
incremento, decremento y complemento a dos mediante la selección de
sus operandos y de \texttt{cin}. La consigna utiliza esta propiedad
para reutilizar una sola instancia de \texttt{adder4} en la unidad
aritmética~\cite{guialab2}. En la resta implementada mediante
$A+\overline{B}+1$, el acarreo de salida indica ausencia de préstamo;
la bandera borrow corresponde a su negación.

\subsection{Diseño jerárquico}
El diseño jerárquico divide un sistema en bloques con interfaces
definidas. Las declaraciones \texttt{component} describen los
puertos de los componentes y las instancias con \texttt{port map}
los conectan mediante señales internas~\cite{jerarquia}.
En una ALU, distintos bloques pueden producir resultados en paralelo
y un multiplexor selecciona la salida correspondiente al
opcode~\cite{casoalu}.

La metodología top-down parte de la función global para definir
bloques e interfaces. La implementación bottom-up comienza por los
bloques básicos y avanza hacia su integración. Esta combinación
permite comprobar cada módulo y localizar errores antes de verificar
el sistema completo. La división en archivos o componentes no
garantiza por sí sola una reducción de recursos o retardos; esos
aspectos requieren evaluar el diseño sintetizado.

\subsection{Testbench y verificación}
Un banco de pruebas instancia el circuito bajo prueba, genera
estímulos y comprueba sus salidas. Las sentencias \texttt{assert}
permiten detectar diferencias respecto de un resultado esperado,
acompañadas de un mensaje que identifique el caso~\cite{fundamentos,casoalu}.
El banco no forma parte del hardware implementado.

La ALU tiene $16\times16\times16=4096$ combinaciones binarias de
opcode, $A$ y $B$. Este espacio permite plantear una comprobación
exhaustiva del resultado y de las banderas. El sumador de cuatro
bits con \texttt{cin} presenta $16\times16\times2=512$
combinaciones. Comparar ambas ALU comprueba su equivalencia, pero
tambi\'en se necesita un modelo de referencia independiente para
detectar errores compartidos por las dos versiones.
